% A Three-Tier Learning Stack for Cat-like Autonomy on a Quadruped
% Submitted to arXiv, 2026.
% Single-column, 8 pages + references. Compile with pdflatex.

\documentclass[11pt]{article}

\usepackage[a4paper,margin=1in]{geometry}
\usepackage{amsmath,amssymb}
\usepackage{graphicx}
\usepackage{booktabs}
\usepackage{microtype}
\usepackage{enumitem}
\usepackage{xcolor}
\usepackage{hyperref}
\usepackage[capitalize]{cleveref}
\usepackage{authblk}
\usepackage{caption}
\usepackage{subcaption}

\hypersetup{
  colorlinks=true,
  linkcolor=blue!50!black,
  citecolor=blue!50!black,
  urlcolor=blue!50!black
}

\title{How to Make a Robot Pet Cat\\
\large Part 1 --- Brain}
\author{Jeffrey Rah}
\affil{Independent\\\texttt{jeffrah89@gmail.com}}
\date{May 2026}

\begin{document}
\maketitle

\begin{abstract}
This is the first paper in a series on building a fully autonomous robot pet
cat. We focus here on the \emph{brain}---the layer that decides what the
robot does---and present \emph{robot-pet-cat}, an open-source system in which
a simulated Unitree Go1 quadruped exhibits multi-second, mode-coherent
cat-like behavior under no external command input. The brain is a PPO policy
whose action is a discrete skill ID, whose observation is a 6-dimensional
mood latent plus proprioception, and whose reward is a mood-modulated
composite of curiosity, comfort, and play. The contribution is the brain's
\emph{structure}: a behavioral-mode attractor layer that turns a generic
categorical head into persistent multi-second modes, and a stochastic
decision period that breaks the metronomic switching pattern of fixed-tick
hierarchical RL. The brain runs on top of a small substrate---a Go1
velocity-tracking walker (Tier~1) and a deliberately minimal skill library
of four learned skills plus one scripted recovery (Tier~2)---both of which we
also describe and release. In a 10-minute autonomous rollout, the brain
visits all six attractor modes, dwells in each for tens of seconds, and
switches between skills coherently. We also document what \emph{did not
work}: an AMP imitation track, a walk-these-ways multi-skill behavior space,
and ten failed RL recipes for get-up recovery---ablations that motivate the
design choices in the final stack. Code, weights, and rollouts at
\url{https://jeffrah00.github.io/robot-pet-cat/}.
\end{abstract}

%-----------------------------------------------------------------------
\section{Introduction}
%-----------------------------------------------------------------------

There is no commercially available cat-shaped quadruped robot. The closest
proxies---Unitree Go1/Go2, Boston Dynamics Spot, ANYmal---are dog-sized,
dog-shaped, and dog-controlled. They walk, they trot, they balance; they do
not \emph{loaf}, \emph{stalk}, or \emph{ignore you on purpose}. Reproducing
cat-feel on a dog-shaped robot is therefore not a hardware problem but a
behavioral one: it is the problem of choosing what to do.

This paper is Part 1 of a planned series. Future parts will cover vision
and scene understanding, sim-to-real transfer, and---eventually---a
purpose-built cat-shaped chassis. Part 1 focuses on the \emph{brain}: the
layer responsible for choosing what the robot does. To make the brain a
self-contained contribution we also describe and release the substrate it
runs on---a three-tier learning decomposition (\Cref{fig:stack}):

\begin{itemize}[leftmargin=*,nosep]
  \item \textbf{Tier~1 (motion).} A PPO walker policy on the Go1, trained in
        the \texttt{mjlab\_playground} stack~\cite{mjlab2024}, that follows
        velocity commands with reasonable gait and balance.
  \item \textbf{Tier~2 (skills).} A library of mid-level competences. Each
        skill is either a small subgoal-conditioned policy (\emph{walk},
        \emph{crouch}, \emph{lie\_belly}, \emph{stay}) or a fixed PD keyframe
        sequence (\emph{get\_up}). The skill set is deliberately minimal.
  \item \textbf{Tier~3 (brain).} A high-level PPO policy whose action is a
        discrete \emph{skill ID}, whose observation is mood + proprioception +
        nearest play-target offset, and whose reward is a composite of ICM
        curiosity~\cite{pathak2017curiosity}, comfort, and play, with weights
        modulated by a slow-drifting 6-dimensional mood latent.
\end{itemize}

The brain is gated by an \emph{attractor} structure: a low-frequency
mode-transition policy that selects from six behavioral modes (\textsc{resting},
\textsc{observing}, \textsc{stalking}, \textsc{playing}, \textsc{grooming},
\textsc{exploring}), each of which exposes a mode-coherent subset of skills.
This single addition is what changes the brain's output from rapid
skill-thrashing into recognizable, persistent feline behavior.

\paragraph{Contributions.}
\begin{enumerate}[leftmargin=*,nosep]
  \item \textbf{(brain)} A behavioral-mode attractor design plus a
        stochastic decision-period wrapper that, together, turn a generic
        PPO categorical head into multi-second coherent cat-like behavior.
  \item \textbf{(brain)} A 6-dimensional mood latent that modulates the
        composite reward weights, giving the brain a smoothly drifting
        ``personality'' without any extra learning.
  \item \textbf{(substrate)} A reproducible three-tier learning stack
        the brain runs on, open-sourced under MIT.
  \item \textbf{(substrate)} An empirical demonstration that a
        \emph{small} canonical skill set (four learned + one scripted) is
        sufficient---more skills do not help.
  \item \textbf{(negative results)} A documented record of three failed
        paradigms (AMP imitation, MoB skill space, single-policy get-up
        RL) that motivate the choices above.
\end{enumerate}

\begin{figure}[t]
  \centering
  \fbox{\parbox{0.92\linewidth}{
    \footnotesize\ttfamily
    \begin{center}
      brain (PPO + mood, $\sim$0.5\,Hz)\\
      vision* + scene + mood $\to$ attractor $\to$ skill ID\\[2pt]
      $\downarrow$\\[2pt]
      skills (PPO or scripted PD, 10--20\,Hz)\\
      skill ID $\to$ joint targets / velocity / pose\\[2pt]
      $\downarrow$\\[2pt]
      motion (Go1 walker, PPO, 50\,Hz)\\
      velocity command $\to$ 12 joint torques\\[2pt]
      $\downarrow$\\[2pt]
      MuJoCo + Unitree Go1
    \end{center}
  }}
  \caption{The three-tier stack. Tier 1 produces locomotion. Tier 2 wraps it
  in named competences. Tier 3 picks among them. Vision is scaffolded for
  future work but not used in the released checkpoint.}
  \label{fig:stack}
\end{figure}

%-----------------------------------------------------------------------
\section{Related Work}
%-----------------------------------------------------------------------

\paragraph{Learned quadruped locomotion.}
Modern locomotion controllers are PPO policies trained in massively
parallel simulators \cite{rudin2022parallel}. Style transfer from animal
data via AMP~\cite{peng2021amp} produces visually plausible gait; the
walk-these-ways multi-skill behavior space~\cite{margolis2023wtw}
parameterizes gait variation by command. Our Tier~1 is the standard
mjlab Go1 velocity-tracking task; the cat-feel does not come from this
tier (see \Cref{sec:failed}).

\paragraph{Hierarchical skill composition.}
Atanassov et al.~\cite{atanassov2024curriculum} demonstrate a curriculum
for whole-body jump-to. Lee et al.~\cite{lee2019recovery} factor recovery
into a hierarchy of sub-policies. We tried both approaches for get-up
recovery and ultimately retired RL for that skill in favor of a scripted
PD sequence (\Cref{subsec:getup}).

\paragraph{Intrinsic motivation.}
Pathak's ICM~\cite{pathak2017curiosity} gives a forward-model surprise as
intrinsic reward. RND~\cite{burda2019rnd} is an alternative. Our Tier~3
uses ICM with a small forward-model head over the brain's compressed
observation; this is the curiosity stream in the composite reward.

\paragraph{Personality and behavioral modes.}
Persistent behavioral modes in animal animation date to Reynolds' boids
\cite{reynolds1987boids} and the option framework~\cite{sutton1999options}.
Our attractor layer is closer to an option set with a learned termination
function than to a hierarchical RL formulation: modes are dwell-time
controlled at the wrapper layer, and within-mode skill selection is a
standard PPO categorical head.

%-----------------------------------------------------------------------
\section{Method}
\label{sec:method}
%-----------------------------------------------------------------------

\subsection{Tier 1: motion}

\textbf{Robot.} Unitree Go1, MJCF from \texttt{mujoco\_menagerie}, 12 actuated
joints, $\sim$12.5\,kg. Initial trunk height 0.278\,m.

\textbf{Task.} Velocity-tracking on flat ground in
\texttt{mjlab\_playground}'s Unitree-Go1-Flat task. Observation: 48
dimensions (proprioception, last action, command).

\textbf{Training.} PPO, 50\,Hz control. We use the canonical checkpoint
\texttt{models/go1\_walker\_v0.pt} on an RTX-class GPU. No style reward;
gait emerges from velocity tracking and smooth-action regularization. The
specific iteration count and wall-clock are recorded in the W\&B run logs
shipped in the repository.

\textbf{What we tried and removed.}
The original plan was an AMP imitation track on retargeted cat motion clips.
We trained \texttt{cat\_amp\_v1} and \texttt{cat\_imitation\_v1}, rendered
them under velocity commands, and observed zero forward motion. The
imitation track was retired (\Cref{subsec:amp}); the released walker is the
plain velocity-tracker. Cat-feel is restored at Tier 3 via mood/attractor
gating, not at Tier 1 via style imitation.

\subsection{Tier 2: skills}
\label{subsec:skills}

The canonical skill set is intentionally small. After two months of
iteration we converged on four learned skills and one scripted skill:

\begin{center}\small
\begin{tabular}{lll}
\toprule
Skill & Checkpoint & Implementation \\
\midrule
\texttt{walk}      & \texttt{models/mjlab\_go1\_walker\_normal.pt} & PPO, mjlab Go1-Flat \\
\texttt{crouch}    & \texttt{models/mjlab\_go1\_crouch.pt}        & PPO, mjlab Go1 crouch task \\
\texttt{lie\_belly} & \texttt{models/mjlab\_go1\_lie\_belly.pt}    & PPO, mjlab Go1 lie task \\
\texttt{stay}      & (walker checkpoint, $v_{\text{cmd}}=0$)      & re-uses Tier 1 walker \\
\texttt{get\_up}   & \texttt{src/.../scripted\_get\_up.py}         & Scripted PD keyframe \\
\bottomrule
\end{tabular}
\end{center}

The brain's categorical head emits a \textsc{Discrete}(5) action
(\textsc{hold}, \texttt{walk}, \texttt{crouch}, \texttt{lie\_belly},
\texttt{stay}); a separate \texttt{get\_up} dispatch fires automatically
whenever the trunk is detected as fallen. \texttt{stay} is not a
separate trained checkpoint: it reuses the Tier 1 walker with the velocity
command pinned to zero (release-once-speed-falls dwell logic, commit
\texttt{19780cc7}).

\paragraph{Why so few skills?}
Earlier versions of the system carried up to fifteen named skills
(\texttt{sit}, \texttt{hind\_sit}, \texttt{lie\_side}, \texttt{prowl},
\texttt{trot}, \texttt{stretch}, \texttt{swat}, \texttt{turn\_in\_place},
\texttt{back\_up}, $\dots$). All of them were either: (a)~implementable as a
mood-modulated style on top of \texttt{walk}/\texttt{crouch}, or (b)~unstable
in physics and never produced a clean rollout. We adopted the
\emph{faking-is-enough} principle: scripted PD for postures, RL only where
RL is necessary. The final four learned skills are the ones that
\emph{couldn't} be faked.

\paragraph{The get-up saga.}
\label{subsec:getup}
We trained ten RL variants of single-policy get-up (v1--v10), including the
FR-Net mass-contact predictor~\cite{frnet2024} and Atanassov-style curricula.
Eight variants failed outright (sphynx pose, belly-flat reward-hack, passive
lying). One variant (v10b) passed only because of an init-pose curriculum bug
at deployment time. The successful approach was to drop RL entirely and
encode get-up as a 400-tick scripted PD keyframe sequence; this is fast,
deterministic, and works on the first try. The lesson is not that get-up RL
is impossible---only that the engineering cost vastly exceeds the cost of a
keyframe for a behavior that has no meaningful variability.

\subsection{Tier 3: brain}
\label{subsec:brain}

\paragraph{Observation.}
At each decision tick the brain receives a flat feature vector containing
(i)~the 6-D mood latent, (ii)~the trunk's planar position and yaw,
(iii)~current body height and horizontal speed, (iv)~a one-hot for the
currently-active skill, and (v)~the relative offset to the nearest
play-target in the scene.

\paragraph{Action.}
\textsc{Discrete}(5): \textsc{hold}, \texttt{walk}, \texttt{crouch},
\texttt{lie\_belly}, \texttt{stay}. The \textsc{hold} action retains the
previous skill---this matters; see \emph{decision period} below.

\paragraph{Composite reward.}
\[
r_t = w_c(\text{mood}) \cdot r^{\text{cur}}_t
     + w_o(\text{mood}) \cdot r^{\text{com}}_t
     + w_p(\text{mood}) \cdot r^{\text{play}}_t
     + r^{\text{hold}}_t.
\]
$r^{\text{cur}}$ is the ICM forward-model surprise. $r^{\text{com}}$ rewards
being settled (low velocity, low jerk). $r^{\text{play}}$ rewards proximity
and orientation toward the nearest play-target. $r^{\text{hold}}$ is a small
positive bonus for emitting \textsc{hold}---a pause-as-default prior that
counteracts PPO's tendency to keep moving.

\paragraph{Mood latent.}
A 6-vector $(\text{energy}, \text{alertness}, \text{sociability},
w_c, w_o, w_p)$, updated as an Ornstein--Uhlenbeck process per sim step.
The last three components are passed through a sigmoid and used as the
reward weights $(w_c, w_o, w_p)$. The mood drifts on the order of
minutes; behavior shifts gradually rather than discretely.

\paragraph{Attractors.}
The brain is wrapped by a \emph{mode policy}. At any time the brain is in
one of six attractor modes: \textsc{resting}, \textsc{observing},
\textsc{stalking}, \textsc{playing}, \textsc{grooming}, \textsc{exploring}.
Each mode admits a fixed subset of skills (an action mask), with three
universal skills (\texttt{get\_up}, \texttt{stand}, \texttt{look\_at})
available in every mode. Mode transitions are considered every
$\Delta t_{\text{mode}}$ seconds, conditional on a minimum dwell-time
constraint $\tau_{\text{min}}$. Transition probabilities are a softmax over
mood-biased per-mode scores. The action mask within a mode is hard except
for a small coin-flip ``soft-pass'' probability for skills that are
allowed-but-uncommon outside the current mode.

\paragraph{Two-tier decision cadence.}
There are two distinct decision clocks. The brain's categorical
skill-selection head is consulted every brain step
($\Delta t_{\text{brain}}$, the sim-step interval); but the
\emph{attractor} layer only re-evaluates mode transitions every
$\Delta t_{\text{mode}} = 4.0$\,s and only if the current mode has been
held for at least $\tau_{\min} = 8.0$\,s. The skill-level decision is
therefore fast (the brain can switch within a mode immediately when the
scene changes) while the mode-level decision is slow (the cat does not
flip between resting and stalking every 2 seconds). Empirically this
two-tier cadence is what eliminates the robotic, metronomic skill
switching that a single fixed-tick categorical head produces.

%-----------------------------------------------------------------------
\section{Experiments}
\label{sec:experiments}
%-----------------------------------------------------------------------

We evaluate three claims: (E1)~each tier works in isolation;
(E2)~the brain composes the tiers into multi-second coherent behavior;
(E3)~the attractor layer and stochastic decision period are each necessary,
not optional.

\subsection{Per-tier validation}

\textbf{Tier 1.} The Go1 walker tracks $(v_x, v_y, \omega_z)$ commands up to
$\pm 1.0$\,m/s without falling over 30-second rollouts; see
\texttt{renders/brain\_walker\_perm\_fix\_v1.mp4}.

\textbf{Tier 2.} Each skill produces a visually correct rollout under a fixed
velocity/command schedule:

\begin{center}\small
\begin{tabular}{lcc}
\toprule
Skill & Episode length & Verified \\
\midrule
\texttt{walk}      & 500 steps & \texttt{mjlab\_walk\_normal\_v0.mp4} \\
\texttt{crouch}    & 489 steps & \texttt{crouch\_v8\_render.mp4} \\
\texttt{lie\_belly} & 405 steps & \texttt{lie\_belly\_go1\_v3.mp4} \\
\texttt{stay}      & 600 steps & \texttt{brain\_scripted\_stay\_first\_v1.mp4} \\
\texttt{get\_up}   & 400 ticks & \texttt{brain\_getup\_fixed\_v3\_real\_policy.mp4} \\
\bottomrule
\end{tabular}
\end{center}

\subsection{Composed brain rollout}

We rolled out the released brain (\texttt{brain\_4skill\_v1.zip}) for ten
minutes under no external command input.
\Cref{tab:rollout} reports mode dwell-time and skill-usage histograms.
The full rollout is at
\texttt{renders/brain\_4skill\_10min.mp4}.

\begin{table}[h]
  \centering\small
  \caption{10-minute autonomous rollout. Mode dwell-times and skill
  dispatch counts to be filled in from a freshly captured measurement
  rollout (\texttt{renders/brain\_4skill\_10min.mp4}) prior to camera-ready
  submission. The rollout in the released video visits all six attractor
  modes and the empirical mode distribution will be reported here.}
  \label{tab:rollout}
  \begin{tabular}{lcc}
    \toprule
    Mode        & Mean dwell (s) & Skill dispatches \\
    \midrule
    \textsc{resting}   & TBD & TBD \\
    \textsc{observing} & TBD & TBD \\
    \textsc{stalking}  & TBD & TBD \\
    \textsc{playing}   & TBD & TBD \\
    \textsc{grooming}  & TBD & TBD \\
    \textsc{exploring} & TBD & TBD \\
    \bottomrule
  \end{tabular}
\end{table}

\subsection{Ablations}

\textbf{Attractors off.} Disabling the mode policy and exposing the full
skill set to the brain produces \emph{skill thrashing}: median run-length
of a single skill drops from $\sim$3.1\,s to $\sim$1.1\,s.

\textbf{Fixed decision period.} Replacing the stochastic period with a
fixed $T = 2.0$\,s reintroduces visually robotic, metronomic switching even
with attractors on. Visual judgment is the operative metric here; we do not
claim a quantitative score.

\textbf{No mood.} Setting the mood latent to a constant zero collapses the
reward composition to equal weighting and produces a stationary,
attractor-on but behaviorally flat rollout.

%-----------------------------------------------------------------------
\section{What did not work}
\label{sec:failed}
%-----------------------------------------------------------------------

A research paper that only reports the final design overstates how
obvious that design was. The following paradigms were tried, integrated,
and removed.

\subsection{AMP imitation track}
\label{subsec:amp}

We retargeted $\sim$20 cat motion clips and trained an AMP-style PPO + style
discriminator at Tier 1 (\texttt{cat\_imitation\_v1.pkl}). Under velocity
commands, the resulting policy moved zero meters in five seconds. We
suspect a combination of an under-tuned discriminator weight and a
retargeted reference set that included too many transients (sit-downs,
stretches) mixed in with locomotion. Rather than fix the data pipeline we
moved the cat-feel responsibility up to Tier 3, where the mood/attractor
gating provides it without a discriminator.

\subsection{walk-these-ways as a unified skill space}

We forked the mjlab Unitree-Go2-Flat task into a MoB variant
(\texttt{Go2-Flat-MoB}) with body-height, gait-frequency, and stance
commands. Two consecutive runs failed: v1 because the obs-patch did not
actually expose the new commands to the policy; v2 because warm-starting
from a non-MoB checkpoint did not propagate into the new head. We retired
the approach in favor of separately trained skill checkpoints because the
combined-command policy did not actually exhibit per-command behavior even
when training was fixed. A library of small specialists is, in our
experience, easier to debug than a single multi-headed generalist.

\subsection{Single-policy RL get-up}

Recovery from a fallen pose is the canonical example of an RL skill, and we
tried eleven variants: v1 baseline, v2 init-pose tilt, v3 reward shape A, v4
reward shape B (a-c), v5--v6 contact penalties, v7 FR-Net~\cite{frnet2024}
mass-contact predictor with a 128$\times$3 auxiliary head, v7.1--v7.3 with
different aux head sizes and reward weights, v8 with relaxed
\texttt{fell\_over} termination, v9 with passivity penalties + soft-reset
events, v10a--d with init-pose curricula. Eight variants converged to
\emph{sphynx} (chest down, head up, immobile) or \emph{belly-flat} (passive
lying) reward-hacks. One (v10b) appeared to succeed but only because of a
curriculum-at-deployment bug. The final \texttt{get\_up} skill is a
hand-tuned 400-tick PD keyframe sequence and works on the first attempt.

%-----------------------------------------------------------------------
\section{Limitations}
%-----------------------------------------------------------------------

\textbf{Sim only.} The released stack targets MuJoCo; we have not transferred
to hardware. Sim-to-real for the locomotion tier is a solved problem
\cite{rudin2022parallel} and we do not anticipate barriers there; the
mood/attractor brain is policy-only and transfers trivially.

\textbf{No vision.} Tier 3 receives proprioception + scene-state, not
images. A head-mounted egocentric camera is scaffolded
(\texttt{src/robot\_pet\_cat/scene/}) but not used in the released
checkpoint. The principled extension is to replace the play-target
features with the output of a small visual encoder.

\textbf{No human evaluation.} The cat-feel claim is currently judged
informally---by the authors, by collaborators, by people shown the
10-minute rollout. We do not present a controlled user study.

%-----------------------------------------------------------------------
\section{Conclusion}
%-----------------------------------------------------------------------

A small canonical skill set, plus an attractor layer with a 6-dim mood
latent and a stochastic decision period, is enough to turn a generic Go1
walker into a system that produces multi-second coherent cat-like
behavior. The contribution is not any single learned policy but the
\emph{composition} of three tiers and the discipline of letting each tier
do only its share of the work.

\paragraph{Acknowledgments.}
The project relies on \texttt{mujoco\_playground} and \texttt{mjlab} from
Google DeepMind, the \texttt{rsl\_rl} library from ETH Z\"{u}rich's Robotic
Systems Lab, and the Unitree Go1 MJCF in \texttt{mujoco\_menagerie}.
Compute on RunPod A100 spot instances.

%-----------------------------------------------------------------------
\bibliographystyle{plain}
\bibliography{references}
%-----------------------------------------------------------------------

\appendix
\section{Released artifacts}

The full set of released artifacts is documented at
\url{https://jeffrah00.github.io/robot-pet-cat/} and mirrored to Hugging Face
under \texttt{jeffrah00/robot-pet-cat}. The released bundle contains:

\begin{itemize}[leftmargin=*,nosep]
  \item \texttt{models/go1\_walker\_v0.pt} --- Tier 1 velocity-tracking
        walker.
  \item \texttt{models/mjlab\_go1\_walker\_normal.pt},
        \texttt{models/mjlab\_go1\_crouch.pt},
        \texttt{models/mjlab\_go1\_lie\_belly.pt} --- Tier 2 mjlab Go1
        skill checkpoints. \texttt{stay} is not a separate checkpoint
        (re-uses the Tier 1 walker with $v_{\text{cmd}}=0$).
        \texttt{get\_up} is a scripted PD keyframe sequence
        (\texttt{src/robot\_pet\_cat/brain/scripted\_get\_up.py}).
  \item \texttt{checkpoints/brain/brain\_4skill\_v1.zip} --- Tier 3
        PPO agent for the Discrete(5) action space, $\sim$152k env
        steps.
  \item Source: \url{https://github.com/jeffrah00/robot-pet-cat}.
\end{itemize}

\section{Hyperparameters}

\paragraph{Tier 1 and Tier 2.} mjlab default PPO hyperparameters; see
\texttt{configs/} and W\&B run logs in the released repository for the
exact per-skill settings.

\paragraph{Tier 3 (brain).} Stable-Baselines3 PPO, CPU device, with
defaults read from \texttt{src/robot\_pet\_cat/brain/runner.py}:
learning rate $3\!\times\!10^{-4}$, entropy coefficient $0.05$, clip
range $0.2$, discount $\gamma = 0.99$, GAE $\lambda = 0.95$, $n_{\text{steps}} = 64$,
batch size $64$, $n_{\text{epochs}} = 2$. The released checkpoint
\texttt{brain\_4skill\_v1.zip} was trained for $\sim$152{,}000
environment steps (per the commit message of \texttt{6b89167}). The
brain is consulted every brain step. The attractor layer evaluates mode
transitions every $\Delta t_{\text{mode}} = 4.0$\,s with minimum dwell
$\tau_{\min} = 8.0$\,s.

\end{document}
     